\scanpage{34}
\mspage{41}
$\underline a=\sum_i m_i(a_i)\in Z(A)$, on posera $\underline q_{\underline a}=\prod_i m_i\,\underline q_{a_i}$. On a toujours $\operatorname{div}(\underline q_{\underline a})=(\operatorname{div}\underline q)_{\underline a}$. Pour $\underline a\in Z_o(A)$, et $v\in M_{\overline K}$, le nombre réel $\underline q_o(\underline a,v)$ ne dépend que de $\underline q$, $\underline a$, et $v$. Nous le désignerons par $\underline q(\underline a,v)$.

\margin{\figure{\begin{tikzpicture}[x=1cm,y=1cm]\draw[gray,line width=.2pt] (0,1.2) -- (1.2,1.65);\draw[gray,line width=.2pt] (0,0) .. controls (.8,-.12) and (1.3,-.16) .. (1.7,-.06);\end{tikzpicture}}}
Soit $X\in D_\ell(V)_{\overline K}$, et soit $f$ une fonction sur $V$, définie sur $\overline K$, telle que $\operatorname{div}(f)=X$. La classe $\chi_V(f)$ de la quasi-fonction principale $\widetilde f$ ne dépend que de $X$. Cette classe sera désignée par $\underset{\sim}{f}(X)$, ou $\underset{\sim}{f}_V(X)$.

Soit $A$ une variété abélienne définie sur $K$. Nous noterons $Q^*_\ell(A)$ (resp. $Q^*_a(A)$) le sous-groupe de $Q^*(A)$ composé des éléments $\underline q$ tels qu'on ait $\operatorname{div}(\underline q)\in D_\ell(A)$ (resp. $\operatorname{div}(\underline q)\in D_a(A)$).

Pour $\underline q\in Q^*_\ell(A)$, nous poserons $\xi_o(\underline q)=\underset{\sim}{f}(\operatorname{div}(\underline q))\,q^{-1}$. On a $\operatorname{div}(\xi_o(\underline q))=0$, d'où $\xi_o(\underline q)\in B^*(A)$, et il est clair que l'application $\xi_o:Q^*_\ell(A)\longrightarrow B^*(A)$ ainsi définie est un homomorphisme. Pour $u\in A_{\overline K}$, l'homomorphisme $\xi_o$ commute avec la translation $\tau_u$; autrement dit, pour $\underline q\in Q^*_\ell(A)$, on a $\xi_o(\underline q_u)=(\xi_o(\underline q))_u$. Montrons en outre que, pour $\underline q\in Q^*(A)$, et pour $\underline a$, $\underline b$ $Z_o(A)$, tels qu'on ait $X_{\underline a}$ et $X_{\underline b^-}\sim0$ (i.e. $\underline q_{\underline a}$ et $\underline q_{\underline b}\in Q^*_\ell(A)$), on a les relations
\[
(15)\qquad (\xi_o(\underline q_{\underline a}))(\underline b,v)=(\xi_o(\underline q_{\underline b^-}))(\underline a^-,v)=(\xi_o(\underline q^-_{\underline b}))(\underline a,v).
\]
En effet, soit $f$ une fonction sur $A$, définie sur $\overline K$, et admettant pour diviseur $X_{\underline a}$. Soit d'autre part $\underline r$ une quasi-fonction sur $A$ représentant $\underline q$. On a
\[
\begin{aligned}
(\xi_o(\underline q_{\underline a}))(\underline b,v)&=(\widetilde f\,r_{\underline a}^{-1})(\underline b,v)\\
&=(\widetilde f\,r_{\underline a}^{-1})_{\underline b^-}(0,v).
\end{aligned}
\]
\mspage{33}
